% ===================================================================
%  ТАБЛИЦА № 10 — Море. — Порт.
%  Traduction 2026 d'apres texto/io, controlee sur texto/fr, texto/en
%  et texto/es. Consigne : voir texto/ru/10-tablica-01.tex.
%
%  Le francais decale sa numerotation d'un cran a partir de l'alinea 8.
%  C'est une coquille de l'imprimeur francais, conservee dans la
%  transcription ; le russe suit l'ido, qui compte juste.
% ===================================================================

%%K t10-apar-1 apar
\VUsaut{4.0mm}
\VUtitre{15.0pt}{60}{ТАБЛИЦА № 10}
\VUsaut{3.0mm}
\VUpk{11.4pt}{40}{Море. --- Порт.}
\VUsaut{3.0mm}

%%K t10-01-1 p
1. --- Летом, пока одни ищут покоя и прохлады в горах, другие
предпочитают ехать на побережье. Так мы и сделали в прошлом году, потому
что очень любим \VUgras{океан} \textsuperscript{(1)}.

%%K t10-02-1 p
2. --- Что до меня, мне доставляет большое удовольствие глядеть на эту
безмерную водную гладь, простёртую до \VUgras{горизонта} \textsuperscript{(2)}, и
смотреть, как огромные \VUgras{пароходы} \textsuperscript{(3)} уходят в долгий
\VUgras{рейс} с путешественниками, направляющимися в Америку.

%%K t10-03-1 p
3. --- Поэтому отец, который знает мои вкусы, всегда выбирает на каникулы
очень тихий берег возле одного из наших больших \VUgras{военных портов}.

%%K t10-03-2 p
В последний наш приезд \VUgras{эскадра} пришла и стала на якорь за
огромным \VUgras{волноломом} \textsuperscript{(4)}, который замыкает и защищает
\VUgras{военную гавань} \textsuperscript{(5)}, и я мог вдоволь налюбоваться
разными \VUgras{военными судами}: от маленькой \VUgras{подводной лодки}
\textsuperscript{(10)}, которая ныряет и исчезает под водой, до огромного
\VUgras{броненосца} \textsuperscript{(9)}, которому до сих пор нечего было
опасаться, кроме нападений \VUgras{миноносца} \textsuperscript{(8)}.

%%K t10-tit-2 sub
\VUsaut{3.0mm}
\VUpk{10.2pt}{40}{Малые суда.}
\VUsaut{3.0mm}

%%K t10-04-1 p
4. --- Этот последний уже умел \VUgras{ловко} находить слабое место в
толстой металлической броне и ставить туда опасную \VUgras{торпеду},
взрыв которой топит судно; но его всё же следовало бояться меньше, чем
подводной лодки, потому что истребители легко пускаются за ним в погоню.

%%K t10-05-1 p
5. --- Я мог видеть и быстрые броненосные \VUgras{крейсеры} \textsuperscript{(6)},
почти столь же неуязвимые, как настоящий броненосец, несколько
\VUgras{транспортов} \textsuperscript{(12)}, три или четыре \VUgras{канонерки}
\textsuperscript{(7)} и, наконец, \VUgras{фрегат} \textsuperscript{(11)}.
\VUgras{Всего занятнее вид этого флота, снимающегося с якоря.}

%%K t10-06-1 p
6. --- Какая разница в сложении между этими громадами стали и изящными
\VUgras{трёхмачтовиками} или грациозными \VUgras{парусниками}
\textsuperscript{(13)}, которые ещё в ходу в торговом флоте и которые я видел
входящими в порт под всеми парусами, гуляя по \VUgras{молу}
\textsuperscript{(14)}. Правда, они много теряли, когда \VUgras{ветер} был против
них. Приходилось убирать все паруса, чтобы войти в док, и нужен был
\VUgras{буксир} \textsuperscript{(16)}, чтобы взять их и подвести к набережной,
как простые \VUgras{баржи} \textsuperscript{(17)}.

%%K t10-tit-3 sub
\VUsaut{3.0mm}
\VUpk{10.2pt}{40}{Торговый порт.}
\VUsaut{3.0mm}

%%K t10-07-1 p
7. --- Порт, о котором я говорю, тем и занятен, что в нём есть не только
военная гавань, созданная искусственно исполинскими работами, но и
чудесный \VUgras{торговый порт} в \VUgras{устье} большой реки.

%%K t10-08-1 p
8. --- Полоса земли, разделяющая оба дока, укреплена и защищена рядом
\VUgras{батарей} и \VUgras{бастионов}, выдвинутых в море. Чтобы гулять по
\VUgras{валам} \textsuperscript{(15)}, нужно особое разрешение. Я легко получил
его через дядю, который служит инженером на \VUgras{верфях}
\textsuperscript{(18)}, и ходил смотреть, как под огромными навесами строят суда.

%%K t10-09-1 p
9. --- Я видел даже спуск броненосца и помню то волнение, какое охватило
всю толпу, когда огромная громада, предоставленная себе, поползла со
своего ложа и всплыла на воде реки.

%%K t10-10-1 p
10. --- Чтобы попасть на верфи, переходишь \VUgras{плац} \textsuperscript{(19)},
где каждый день учатся войска гарнизона, и идёшь по железному
\VUgras{мостику} \textsuperscript{(20)}, соединяющему плац с \VUgras{арсеналом}
\textsuperscript{(21)}.

%%K t10-tit-4 sub
\VUsaut{3.0mm}
\VUpk{10.2pt}{40}{Арсенал и доки.}
\VUsaut{3.0mm}

%%K t10-11-1 p
11. --- С дядей я осмотрел это большое заведение, полное \VUgras{пушек}
\textsuperscript{(22)}, \VUgras{ядер} \textsuperscript{(23)} и \VUgras{снарядов}. Видел
я и \VUgras{котельную} \textsuperscript{(24)}, где делают \VUgras{котлы}
\textsuperscript{(25)} для пароходов.

%%K t10-12-1 p
12. --- Совсем рядом --- \VUgras{доки} \textsuperscript{(26)}, сухие доки, и очень
примечательные \VUgras{плавучие доки} \textsuperscript{(27)}, в которых я мог
вблизи разглядеть на чинящемся судне \VUgras{винт} \textsuperscript{(28)},
\VUgras{киль} \textsuperscript{(29)} и \VUgras{руль} \textsuperscript{(30)} корабля.

%%K t10-13-1 p
13. --- Торговый порт стоит изучения не меньше военного, и я провёл
немало часов, слоняясь по набережным, где стоят \VUgras{колёсные пароходы}
\textsuperscript{(31)}, которые ходят вверх по реке, и глядя, как работает
\VUgras{шеровый кран} \textsuperscript{(32)}, поднимая огромные тяжести, точно
пёрышки.

%%K t10-tit-5 sub
\VUsaut{4.0mm}
\VUpk{10.2pt}{40}{Лоцман.}
\VUsaut{3.0mm}

%%K t10-14-1 p
14. --- Я любовался \VUgras{трансатлантическими пароходами}
\textsuperscript{(34)}, покидающими рейд с развевающимися \VUgras{флагами}
\textsuperscript{(35)}, которых вёл искусный \VUgras{лоцман} \textsuperscript{(36)},
дивно умевший держаться \VUgras{фарватера}, размеченного \VUgras{буями}
\textsuperscript{(40)} и \VUgras{вехами}, обходя \VUgras{лихтеры}
\textsuperscript{(41)}, которыми так трудно править под единственным
четырёхугольным парусом. Я различал на \VUgras{носу} \textsuperscript{(37)}
\VUgras{каюты} \textsuperscript{(39)} второго класса, а на \VUgras{корме}
\textsuperscript{(38)} --- салоны для пассажиров первого.

%%K t10-15-1 p
15. --- Иногда я смотрел и на погрузку \VUgras{грузового судна}
\textsuperscript{(42)}, в которое сваливали товары со \VUgras{склада}
\textsuperscript{(43)} и с \VUgras{портовой станции} \textsuperscript{(44)},
подвозимые многочисленными поездами \VUgras{набережной ветки}
\textsuperscript{(45)}.

%%K t10-tit-6 sub
\VUsaut{3.0mm}
\VUpk{10.2pt}{40}{Катанье на лодке.}
\VUsaut{3.0mm}

%%K t10-16-1 p
16. --- Я часто катался на лодке в \VUgras{закрытом доке} \textsuperscript{(53)},
где укрываются маленькие \VUgras{лодки} \textsuperscript{(46)}, и для меня было
радостью держать \VUgras{весло} \textsuperscript{(47)}. Я поднимался на
\VUgras{палубу} \textsuperscript{(48)} \VUgras{парусной лодки} \textsuperscript{(49)},
влезал по \VUgras{снастям} \textsuperscript{(52)} на \VUgras{мачту}
\textsuperscript{(51)} до \VUgras{рей} \textsuperscript{(50)}, а потом продолжал
прогулку в \VUgras{яле} \textsuperscript{(54)} среди тяжёлых \VUgras{рыбачьих
судов} \textsuperscript{(55)}, где сохнут развешанные \VUgras{сети}
\textsuperscript{(56)}.

%%K t10-17-1 p
17. --- На \VUgras{набережной} \textsuperscript{(58)} передо мною проходили самые
разные \VUgras{товары} \textsuperscript{(59)}, упакованные в \VUgras{ящики}
\textsuperscript{(60)} или в \VUgras{тюки} \textsuperscript{(61)}. Они составляли
\VUgras{груз} \textsuperscript{(63)} барж, и мощные \VUgras{краны}
\textsuperscript{(62)} вынимали их и ставили на \VUgras{грузовики}
\textsuperscript{(64)}, пока \VUgras{перевозчики} объявляли их и платили пошлину
в \VUgras{таможне} \textsuperscript{(65)}.

%%K t10-18-1 p
18. --- Наконец, дойдя до \VUgras{разводного моста} \textsuperscript{(66)} или до
\VUgras{шлюза} \textsuperscript{(69)}, миновав \VUgras{землечерпалку}
\textsuperscript{(68)}, которая чистит \VUgras{канал} \textsuperscript{(67)}, я
приставал к молу возле \VUgras{сигнальной мачты} \textsuperscript{(70)}.

%%K t10-19-1 p
19. --- Именно с другой стороны этого мола пристают \VUgras{катера}
\textsuperscript{(71)}, привозящие на берег офицеров эскадры. Красиво смотреть,
как матросы разом поднимают вёсла, а шлюпка, несомая \VUgras{волной}
\textsuperscript{(72)}, легко подходит к сходням. Отсюда лучше всего видны
\VUgras{прилив} и движение \VUgras{отлива} и \VUgras{прибоя} на
\VUgras{пляже} \textsuperscript{(73)}.

%%K t10-tit-7 sub
\VUsaut{3.0mm}
\VUpk{10.2pt}{40}{Офицерское собрание.}
\VUsaut{3.0mm}

%%K t10-20-1 p
20. --- Чтобы вернуться домой, я садился в \VUgras{электрический трамвай}
\textsuperscript{(74)}, \VUgras{остановка} \textsuperscript{(75)} которого у
\VUgras{столба} перед офицерским собранием.

%%K t10-20-2 p
С \VUgras{балкона} \textsuperscript{(76)} собрания, украшенного \VUgras{якорями}
\textsuperscript{(77)}, пушками и ядрами, я часто видел блестящее общество
морских офицеров: \VUgras{лейтенантов} \textsuperscript{(78)} со шпагами,
\VUgras{капитанов артиллерии} \textsuperscript{(79)} и \VUgras{пехоты}
\textsuperscript{(80)} с \VUgras{саблями}. Позади них стоял \VUgras{матрос}
\textsuperscript{(81)}, который подавал им \VUgras{подзорную трубу}, когда они
хотели рассмотреть, что делается вдали.

%%K t10-21-1 p
21. --- Видел я и молодых \VUgras{мичманов} \textsuperscript{(82)}, получавших
приказания от своего \VUgras{командира} \textsuperscript{(83)} или от
\VUgras{адмирала} \textsuperscript{(84)}, пока \VUgras{квартирмейстер}
\textsuperscript{(85)} ждал, чтобы отвезти их на судно.

%%K t10-22-1 p
22. --- С этого балкона вид великолепен: он господствует над всем пляжем
с его \VUgras{палатками} \textsuperscript{(86)} и \VUgras{купальными кабинами}
\textsuperscript{(87)}, и в час купанья видно, как \VUgras{купальщики}
\textsuperscript{(88)} весело резвятся перед \VUgras{казино} \textsuperscript{(89)}.

%%K t10-23-1 p
23. --- В конце пляжа берег образует маленький скалистый \VUgras{полуостров}
\textsuperscript{(91)}, о который разбиваются \VUgras{буруны} \textsuperscript{(92)} и
на котором построен \VUgras{маяк} \textsuperscript{(93)}, ночью указывающий
морякам путь. Этот полуостров соединён с землёй лишь очень узким
\VUgras{перешейком} \textsuperscript{(90)}.

%%K t10-tit-8 sub
\VUsaut{3.0mm}
\VUpk{10.2pt}{40}{Общий вид берега.}
\VUsaut{3.0mm}

%%K t10-24-1 p
24. --- \VUgras{Утёс} \textsuperscript{(94)} тянется дальше, изрезанный морем,
которое вытачивает в нём причудливую череду \VUgras{бухт} \textsuperscript{(95)}
и \VUgras{мысов}, делая его очень живописным.

%%K t10-24-2 p
На \VUgras{плоскогорье} \textsuperscript{(96)}, господствующем над \VUgras{проливом}
\textsuperscript{(98)}, стоят очень важный \VUgras{форт} \textsuperscript{(97)} и
несколько вилл.

%%K t10-25-1 p
25. --- Этот пролив отделяет от материка очень опасный \VUgras{остров}
\textsuperscript{(99)}, окружённый \VUgras{скалами} \textsuperscript{(100)} и
\VUgras{рифами}, где иногда гибнут лодки и даже большие суда. Как ни
проворна помощь и как ни быстро приходят \VUgras{спасательные шлюпки}, эти
\VUgras{кораблекрушения} ужасны, и в равноденственные \VUgras{бури}
несколько судов погибло со всем экипажем и грузом.

%%K t10-25-2 p
Несмотря на такое соседство, \VUgras{порт} этот моряки посещают охотно.
\VUsaut{6.0mm}
\VUfilet{22mm}
